% ╔══════════════════════════════════════════════════════════════════════════╗
% ║  ESCUELA COLOMBIANA DE INGENIERÍA JULIO GARAVITO                        ║
% ║  Software Architecture (ARSW) -- Laboratory No. 8                       ║
% ║  Terraform on Azure -- Load Balancing & Infrastructure as Code          ║
% ╚══════════════════════════════════════════════════════════════════════════╝

\documentclass[10pt,a4paper,twocolumn]{article}

% --- Encoding & language ---------------------------------------------------
\usepackage[utf8]{inputenc}
\usepackage[english]{babel}
\usepackage{lmodern}      % Scalable Latin Modern fonts (microtype expansion)
\usepackage[T1]{fontenc}

% --- Maths -----------------------------------------------------------------
\usepackage{amsmath}

% --- Graphics & layout -----------------------------------------------------
\usepackage{graphicx}
\usepackage[top=2cm, bottom=2cm, left=1.8cm, right=1.8cm]{geometry}
\usepackage{fancyhdr}
\usepackage{float}
\usepackage{caption}
\usepackage{subcaption}

% --- Code listing ----------------------------------------------------------
\usepackage{listings}

% --- Colors & boxes --------------------------------------------------------
\usepackage[dvipsnames]{xcolor}
\usepackage[most]{tcolorbox}

% --- Hyperlinks ------------------------------------------------------------
\usepackage{hyperref}

% --- Typography helpers ----------------------------------------------------
\usepackage{enumitem}
\usepackage{booktabs}
\usepackage{titlesec}
\usepackage{microtype}
\usepackage{parskip}
\usepackage{array}
\usepackage{makecell}

% ===========================================================================
%  COLOUR & STYLE DEFINITIONS
% ===========================================================================
\definecolor{darkblue}{RGB}{0,32,96}
\definecolor{azureblue}{RGB}{0,120,212}
\definecolor{terraformpurple}{RGB}{123,66,188}
\definecolor{lightblue}{RGB}{173,216,230}
\definecolor{codebg}{RGB}{245,245,245}
\definecolor{darkgreen}{RGB}{0,128,0}
\definecolor{placeholderbg}{RGB}{240,240,240}
\definecolor{successgreen}{RGB}{40,167,69}
\definecolor{warningorange}{RGB}{255,140,0}
\definecolor{dangerred}{RGB}{220,53,69}

% Section header color
\titleformat{\section}{\large\bfseries\color{darkblue}}{}{0em}{\thesection\quad}
\titleformat{\subsection}{\normalsize\bfseries\color{darkblue!80}}{}{0em}{\thesubsection\quad}
\titlespacing*{\section}{0pt}{10pt}{4pt}
\titlespacing*{\subsection}{0pt}{8pt}{3pt}

% --- Header / Footer -------------------------------------------------------
\pagestyle{fancy}
\fancyhf{}
\lhead{\textcolor{darkblue}{\textbf{Software Architecture}}}
\rhead{\textcolor{darkblue}{Lab 08 -- Terraform Azure Load Balancing}}
\cfoot{\thepage}
\renewcommand{\headrulewidth}{0.4pt}

% --- Code Listing Styles ---------------------------------------------------
\lstdefinestyle{hclcode}{
    backgroundcolor=\color{codebg},
    basicstyle=\ttfamily\scriptsize,
    keywordstyle=\color{terraformpurple}\bfseries,
    commentstyle=\color{gray}\itshape,
    stringstyle=\color{darkgreen},
    numberstyle=\tiny\color{gray},
    numbers=left,
    stepnumber=1,
    numbersep=5pt,
    frame=single,
    framerule=0.4pt,
    rulecolor=\color{gray!50},
    breaklines=true,
    showstringspaces=false,
    tabsize=2,
    captionpos=b,
    xleftmargin=10pt,
    xrightmargin=5pt,
    aboveskip=6pt,
    belowskip=6pt,
    morekeywords={resource,module,variable,output,provider,terraform,
                  required_providers,required_version,backend,locals,
                  data,for,in,if,else,true,false,null,for_each,count,
                  source,version,type,default,description,sensitive,
                  validation,condition,error_message},
    morecomment=[l]{\#}
}
\lstdefinestyle{yamlcode}{
    backgroundcolor=\color{codebg},
    basicstyle=\ttfamily\scriptsize,
    keywordstyle=\color{azureblue}\bfseries,
    commentstyle=\color{gray}\itshape,
    stringstyle=\color{Brown},
    frame=single, framerule=0.4pt, rulecolor=\color{gray!50},
    breaklines=true, showstringspaces=false, tabsize=2,
    captionpos=b, xleftmargin=10pt, aboveskip=6pt, belowskip=6pt,
    morecomment=[l]{\#}
}
\lstdefinestyle{shellcode}{
    backgroundcolor=\color{black!90},
    basicstyle=\ttfamily\scriptsize\color{white},
    keywordstyle=\color{cyan}\bfseries,
    commentstyle=\color{gray}\itshape,
    stringstyle=\color{yellow},
    frame=single, framerule=0.4pt, rulecolor=\color{gray!50},
    breaklines=true, showstringspaces=false, tabsize=2,
    captionpos=b, xleftmargin=10pt, aboveskip=6pt, belowskip=6pt,
    morekeywords={az,terraform,curl,export,echo,cd,ssh,ssh-keygen,git,gh},
    morecomment=[l]{\#}
}

% --- Custom tcolorbox environments -----------------------------------------
\tcbset{
    conceptbox/.style={
        colback=blue!5!white, colframe=darkblue,
        coltitle=white, colbacktitle=darkblue,
        title={\textbf{Concept}}, fonttitle=\bfseries\small,
        boxrule=0.8pt, before skip=8pt, after skip=8pt,
        left=4pt, right=4pt, top=6pt, bottom=6pt, breakable
    },
    resultbox/.style={
        colback=green!5!white, colframe=successgreen,
        coltitle=white, colbacktitle=successgreen,
        title={\textbf{Result}}, fonttitle=\bfseries\small,
        boxrule=0.8pt, before skip=8pt, after skip=8pt,
        left=4pt, right=4pt, top=6pt, bottom=6pt, breakable
    },
    note/.style={
        colback=yellow!10!white, colframe=warningorange,
        coltitle=white, colbacktitle=warningorange,
        title={\textbf{Note}}, fonttitle=\bfseries\small,
        boxrule=0.8pt, before skip=8pt, after skip=8pt,
        left=4pt, right=4pt, top=6pt, bottom=6pt, breakable
    },
    warningbox/.style={
        colback=red!5!white, colframe=dangerred,
        colbacktitle=dangerred, coltitle=white,
        title={\textbf{Key Design Decision}}, fonttitle=\bfseries\small,
        boxrule=0.8pt, before skip=8pt, after skip=8pt,
        left=4pt, right=4pt, top=6pt, bottom=6pt, breakable
    }
}
\newtcolorbox{conceptbox}[1][]{conceptbox, #1}
\newtcolorbox{resultbox}[1][]{resultbox, #1}
\newtcolorbox{note}[1][]{note, #1}
\newtcolorbox{warningbox}[1][]{warningbox, #1}

% --- Image helper ----------------------------------------------------------
\newcommand{\imgplaceholder}[2]{%
    \IfFileExists{#2}{%
        \includegraphics[height=#1,width=\linewidth,keepaspectratio]{#2}%
    }{%
        \begin{tcolorbox}[
            colback=placeholderbg, colframe=gray!60,
            height=#1, valign=center, halign=center,
            fontupper=\small\ttfamily\color{gray!70},
            boxrule=0.5pt, left=2pt, right=2pt
        ]
            #2
        \end{tcolorbox}%
    }%
}

\setlist[itemize]{leftmargin=15pt, itemsep=2pt, parsep=0pt, topsep=5pt}
\setlist[enumerate]{leftmargin=15pt, itemsep=2pt, parsep=0pt, topsep=5pt}
\captionsetup{font=small, labelfont=bf, format=hang, indention=0pt, margin=10pt}
\hypersetup{colorlinks=true, linkcolor=darkblue, urlcolor=darkblue,
            citecolor=darkblue, pdfborder={0 0 0}}

% ----------------------- TITLE PAGE ----------------------------------------
\title{\vspace{-1cm}
    \begin{center}
        \imgplaceholder{3cm}{media/university_logo.png}
    \end{center}
    \vspace{1.5cm}
    \textbf{\Large Software Architecture Laboratory}\\
    \vspace{1cm}
    \textbf{\huge Laboratory No. 8}\\
    \textbf{\huge Infrastructure as Code with Terraform on Azure}\\
    \vspace{1.5cm}
    \large Public Load Balancer / Linux VMs / Remote State / GitHub Actions OIDC
}

\author{
    \vspace{2cm}
    \textbf{Students:}\\[0.05cm]
    Andersson David S\'{a}nchez M\'{e}ndez\\[0.3cm]
    Cristian Santiago Pedraza Rodr\'{i}guez\\[0.3cm]
    Elizabeth Correa Su\'{a}rez\\[0.3cm]
    Juan Sebastian Ortega Mu\~{n}oz\\[1.5cm]
    \textbf{Instructor:} Professor Javier Ivan Toquica Barrera\\[0.5cm]
    \textbf{Course:} Software Architecture (ARSW)\\[0.3cm]
    \textbf{Institution:} Escuela Colombiana de Ingenier\'{i}a Julio Garavito\\[2cm]
}

\date{May 2026}

% ===========================================================================
\begin{document}

\onecolumn
\maketitle
\thispagestyle{empty}
\newpage

\tableofcontents
\newpage

% ===========================================================================
\twocolumn
% ===========================================================================

% ---------------------------------------------------------------------------
\section{Objective}
% ---------------------------------------------------------------------------

This laboratory modernises the classic Azure load-balancing exercise by
expressing the entire infrastructure as code with \textbf{Terraform}.
Building on the cloud foundations established in previous laboratories,
the team designed and provisioned a reproducible deployment that exposes
two Linux web servers behind a public, Layer-4 \textbf{Azure Load Balancer},
backs the configuration with a remote, lock-protected Terraform
\textbf{state}, and integrates a \textbf{GitHub Actions} pipeline that
authenticates against Azure using \textbf{OpenID Connect} federation
instead of long-lived secrets.

\begin{conceptbox}[title={Lab \#8 Objectives}]
\begin{enumerate}
  \item Model an Azure deployment with Terraform: providers, modules,
        variables, outputs and remote state.
  \item Provision a high-availability workload behind a Standard SKU
        Load Balancer with at least two Linux backends.
  \item Apply a minimum security baseline: NSG with explicit allow-list,
        SSH key-only authentication, naming conventions and tags.
  \item Externalise state to an Azure Storage Account using blob lease
        locking to prevent concurrent applies.
  \item Build a CI/CD pipeline in GitHub Actions that performs
        \texttt{plan} on every pull request and gated \texttt{apply} on
        manual dispatch, all under OIDC.
  \item Validate that the LB round-robins traffic across the backends,
        observe approximate costs and destroy the deployment safely.
\end{enumerate}
\end{conceptbox}

% ---------------------------------------------------------------------------
\section{Target Architecture}
% ---------------------------------------------------------------------------

The deployment is partitioned into three Terraform modules so that each
concern is independently versionable and reviewable. A single Resource
Group acts as the container, which keeps the cleanup operation atomic.

\begin{center}
\small
\begin{tabular}{@{}ll@{}}
\toprule
\textbf{Module} & \textbf{Responsibility} \\
\midrule
\texttt{vnet}    & Virtual Network and the \texttt{web} / \texttt{mgmt} subnets. \\
\texttt{compute} & NICs and Linux VMs bootstrapped by cloud-init. \\
\texttt{lb}      & Public Load Balancer, health probe, rule and NSG. \\
\bottomrule
\end{tabular}
\end{center}

\begin{warningbox}
The choice of \textbf{Layer-4 Load Balancer} over Application Gateway
(L7) is deliberate. The workload serves identical static HTML over plain
HTTP, with no path routing, no TLS termination and no WAF requirements.
A Standard LB delivers sub-millisecond latency and one-tenth of the
hourly cost of Application Gateway. Should the workload need TLS or
path-based routing in the future, only the \texttt{lb} module would be
swapped, leaving \texttt{vnet} and \texttt{compute} untouched. This is
the precise benefit modular Terraform was designed to deliver.
\end{warningbox}

\subsection{Component diagram}

The component diagram (Figure~\ref{fig:component}) is rendered with
\textbf{advanced Mermaid notation} (subgraphs, classed nodes, custom
theming) and captures the four logical layers of the deployment:
identity (Azure AD + OIDC), remote state (Storage Account + container
+ blob lease), workload (Resource Group containing VNet, NSG, NICs,
VMs and the LB) and the external actors (developer CLI, GitHub runner
and the end user). The full Mermaid source lives in
\texttt{docs/DIAGRAMS.md}.

\begin{figure}[H]
\centering
\imgplaceholder{7cm}{media/component-diagram.png}
\caption{Component diagram of the Lab \#8 deployment: identity
federation, remote state backend, edge layer and load-balanced VMs.}
\label{fig:component}
\end{figure}

\subsection{Sequence diagram --- end-user request}

Figure~\ref{fig:seq-enduser} depicts what happens when an end user
issues \texttt{curl http://<lb\_public\_ip>}: name resolution, TCP
handshake intercepted by the LB, 5-tuple hashing that selects a
backend NIC, NSG acceptance of the inbound \texttt{80/TCP} flow, and
the response returning back through the same path. The diagram also
visualises the periodic TCP/80 health probes that keep the backend
pool membership accurate.

\begin{figure}[H]
\centering
\imgplaceholder{7cm}{media/sequence-diagram-end-user.png}
\caption{Request flow through the Public IP, the Load Balancer
5-tuple hash and the NSG before the nginx VM responds with the page
that embeds its hostname.}
\label{fig:seq-enduser}
\end{figure}

\subsection{Sequence diagram --- CI/CD with OIDC}

Figure~\ref{fig:seq-cicd} complements the previous one by showing the
operational path: an engineer triggers a workflow, GitHub mints a
short-lived OIDC JWT, Azure AD swaps it for an access token, the
runner acquires a blob lease on the state file, runs Terraform against
ARM and finally posts the plan back to the PR.

\begin{figure}[H]
\centering
\imgplaceholder{7cm}{media/sequence-diagram-ci-cd.png}
\caption{CI/CD sequence: GitHub Actions trades the OIDC JWT for an
Azure AD token, locks the remote state and runs Terraform.}
\label{fig:seq-cicd}
\end{figure}

% ---------------------------------------------------------------------------
\section{Part 1 --- Modular Terraform Code}
% ---------------------------------------------------------------------------

The root composition wires the three modules together while keeping
provider configuration centralised in \texttt{providers.tf}. The
\texttt{azurerm} backend block is intentionally empty, so the actual
storage configuration is injected at \texttt{terraform init} time via
\texttt{-backend-config=backend.hcl}, keeping subscription identifiers
out of source control.

\begin{lstlisting}[style=hclcode, caption={Root composition (\texttt{infra/main.tf}) wiring the three modules}]
resource "azurerm_resource_group" "rg" {
  name     = "${var.prefix}-rg"
  location = var.location
  tags     = var.tags
}

module "vnet" {
  source              = "../modules/vnet"
  resource_group_name = azurerm_resource_group.rg.name
  location            = azurerm_resource_group.rg.location
  prefix              = var.prefix
  tags                = var.tags
}

module "compute" {
  source              = "../modules/compute"
  resource_group_name = azurerm_resource_group.rg.name
  location            = azurerm_resource_group.rg.location
  prefix              = var.prefix
  admin_username      = var.admin_username
  ssh_public_key      = file(var.ssh_public_key)
  subnet_id           = module.vnet.subnet_web_id
  vm_count            = var.vm_count
  cloud_init          = file("${path.module}/cloud-init.yaml")
  tags                = var.tags
}

module "lb" {
  source              = "../modules/lb"
  resource_group_name = azurerm_resource_group.rg.name
  location            = azurerm_resource_group.rg.location
  prefix              = var.prefix
  backend_nic_ids     = module.compute.nic_ids
  allow_ssh_from_cidr = var.allow_ssh_from_cidr
  tags                = var.tags
}
\end{lstlisting}

\subsection{Variable validation}

Every root variable is typed and validated where it makes sense, so the
plan fails fast on bad inputs instead of producing partial deployments.
For example, \texttt{vm\_count} is bounded between 2 and 5 to enforce
high availability without letting the lab spend run away.

\begin{lstlisting}[style=hclcode, caption={Validated input variables (excerpt of \texttt{infra/variables.tf})}]
variable "prefix" {
  type        = string
  description = "Short identifier prepended to every resource name."
  validation {
    condition     = can(regex("^[a-z][a-z0-9-]{1,10}$", var.prefix))
    error_message = "prefix must be 2-11 lowercase alphanumerics."
  }
}

variable "vm_count" {
  type        = number
  description = "Number of Linux VMs registered behind the LB."
  validation {
    condition     = var.vm_count >= 2 && var.vm_count <= 5
    error_message = "vm_count must be between 2 and 5."
  }
}
\end{lstlisting}

\subsection{Cloud-init bootstrap}

The Linux VMs install nginx on first boot and write a tiny banner that
embeds the hostname. This makes the LB round-robin behaviour observable
end-to-end with a simple \texttt{curl}.

\begin{lstlisting}[style=yamlcode, caption={\texttt{infra/cloud-init.yaml}: nginx installation with hostname banner}]
#cloud-config
package_update: true
packages:
  - nginx
runcmd:
  - echo "<h1>Hello from $(hostname)</h1>" > /var/www/html/index.nginx-debian.html
  - systemctl enable nginx
  - systemctl restart nginx
\end{lstlisting}

% ---------------------------------------------------------------------------
\section{Part 2 --- Load Balancer \& NSG Module}
% ---------------------------------------------------------------------------

The \texttt{lb} module is where the design tradeoffs concentrate. It
provisions a Standard SKU Public IP, a Standard SKU Load Balancer with a
single public frontend, the backend address pool, the TCP/80 health
probe, the load balancing rule and the NSG. The NSG is bound at NIC
level (rather than subnet level) so subsequent tiers can declare their
own posture without inheritance surprises.

\begin{lstlisting}[style=hclcode, caption={Standard LB, probe, rule and NSG (excerpt of \texttt{modules/lb/main.tf})}]
resource "azurerm_lb_probe" "probe" {
  name            = "http-80"
  loadbalancer_id = azurerm_lb.lb.id
  protocol        = "Tcp"
  port            = 80
}

resource "azurerm_lb_rule" "rule" {
  name                           = "http-80"
  loadbalancer_id                = azurerm_lb.lb.id
  protocol                       = "Tcp"
  frontend_port                  = 80
  backend_port                   = 80
  frontend_ip_configuration_name = "Public"
  backend_address_pool_ids       = [azurerm_lb_backend_address_pool.bepool.id]
  probe_id                       = azurerm_lb_probe.probe.id
}

resource "azurerm_network_security_group" "nsg" {
  name                = "${var.prefix}-web-nsg"
  location            = var.location
  resource_group_name = var.resource_group_name

  security_rule {
    name                       = "Allow-HTTP-Internet"
    priority                   = 100
    direction                  = "Inbound"
    access                     = "Allow"
    protocol                   = "Tcp"
    destination_port_range     = "80"
    source_address_prefix      = "*"
    source_port_range          = "*"
    destination_address_prefix = "*"
  }

  security_rule {
    name                       = "Allow-SSH-From-Operator"
    priority                   = 110
    direction                  = "Inbound"
    access                     = "Allow"
    protocol                   = "Tcp"
    destination_port_range     = "22"
    source_address_prefix      = var.allow_ssh_from_cidr
    source_port_range          = "*"
    destination_address_prefix = "*"
  }
}
\end{lstlisting}

\begin{note}
The Standard SKU was selected over Basic because Public IPs of Standard
SKU are mandatory for the Standard LB family, and Basic SKU is on the
deprecation track in Azure since September 2025. Choosing Standard
preserves availability zone readiness without committing to it on day
one.
\end{note}

% ---------------------------------------------------------------------------
\section{Part 3 --- Remote State \& Locking}
% ---------------------------------------------------------------------------

Local Terraform state is incompatible with collaborative work: any
parallel \texttt{apply} would race and corrupt the file. The team
provisioned a dedicated Storage Account
(\texttt{sttfstate658} in resource group \texttt{rg-tfstate-lab8}) with
a single private container (\texttt{tfstate}). Terraform automatically
acquires a blob lease on the state file before each operation, which
guarantees mutual exclusion and survives crashed runs (the lease auto-
expires after one minute).

\begin{lstlisting}[style=shellcode, caption={Bootstrap of the remote state backend (run once)}]
SUFFIX=$RANDOM
LOCATION=eastus
RG=rg-tfstate-lab8
STO=sttfstate${SUFFIX}
CONTAINER=tfstate

az group create -n $RG -l $LOCATION
az storage account create -g $RG -n $STO -l $LOCATION \
   --sku Standard_LRS --encryption-services blob
az storage container create --name $CONTAINER --account-name $STO
\end{lstlisting}

The corresponding \texttt{backend.hcl} (kept out of source control) is
the only file containing storage names; \texttt{providers.tf} declares
an empty \texttt{backend "azurerm" \{\}} block that the CLI hydrates at
init time.

% ---------------------------------------------------------------------------
\section{Part 4 --- GitHub Actions with OIDC}
% ---------------------------------------------------------------------------

Storing an Azure Service Principal client secret inside GitHub
\texttt{Secrets} works, but rotates poorly and exposes a long-lived
credential on every job log leak. Instead, the team configured
\textbf{OpenID Connect federation} between GitHub Actions and Azure AD.
Each workflow run mints a short-lived OIDC JWT, swaps it for an Azure
AD access token via \texttt{azure/login@v2}, and uses that ephemeral
token to call ARM and the state Storage Account. The federation trust
itself lives entirely on the Azure side as an \emph{App Registration
federated credential} bound to the GitHub repository's OIDC subject.

\begin{lstlisting}[style=yamlcode, caption={Lint + Plan jobs (excerpt of \texttt{.github/workflows/terraform.yml})}]
permissions:
  id-token: write       # required for OIDC token request
  contents: read
  pull-requests: write

jobs:
  lint:
    runs-on: ubuntu-latest
    steps:
      - uses: actions/checkout@v4
      - uses: hashicorp/setup-terraform@v3
        with: { terraform_version: "1.9.5" }
      - run: terraform fmt -check -recursive
      - uses: azure/login@v2
        with:
          client-id:       ${{ secrets.AZURE_CLIENT_ID }}
          tenant-id:       ${{ secrets.AZURE_TENANT_ID }}
          subscription-id: ${{ secrets.AZURE_SUBSCRIPTION_ID }}
      - run: terraform init -backend=false
        working-directory: infra
      - run: terraform validate -no-color
        working-directory: infra

  plan:
    needs: lint
    if: github.event_name == 'pull_request'
    runs-on: ubuntu-latest
    steps:
      - uses: actions/checkout@v4
      - uses: hashicorp/setup-terraform@v3
        with: { terraform_version: "1.9.5" }
      - uses: azure/login@v2
        with:
          client-id:       ${{ secrets.AZURE_CLIENT_ID }}
          tenant-id:       ${{ secrets.AZURE_TENANT_ID }}
          subscription-id: ${{ secrets.AZURE_SUBSCRIPTION_ID }}
      - working-directory: infra
        run: |
          terraform init \
            -backend-config="resource_group_name=${{ secrets.BACKEND_RG }}" \
            -backend-config="storage_account_name=${{ secrets.BACKEND_SA }}" \
            -backend-config="container_name=${{ secrets.BACKEND_CONTAINER }}" \
            -backend-config="key=${{ secrets.BACKEND_KEY }}"
          terraform plan -var-file=env/dev.tfvars -no-color -out=tfplan
\end{lstlisting}

\begin{resultbox}[title={Pipeline behaviour}]
\begin{itemize}
  \item Every pull request to \texttt{main} runs \texttt{fmt},
        \texttt{validate} and \texttt{plan}; the plan diff is posted
        back to the PR as a sticky comment so reviewers see the impact
        without having to fetch the branch locally.
  \item The \texttt{apply} and \texttt{destroy} jobs are gated behind
        the \texttt{production} GitHub Environment, which requires
        explicit reviewer approval and restricts deployments to the
        \texttt{main} branch.
  \item No long-lived secret leaves Azure: the \texttt{azure/login@v2}
        step exchanges the GitHub OIDC JWT for a 1-hour AAD token.
\end{itemize}
\end{resultbox}

% ---------------------------------------------------------------------------
\section{Part 5 --- Operational Validation}
% ---------------------------------------------------------------------------

After a successful \texttt{terraform apply}, the team verified the
deployment end-to-end. Outputs of the recorded run:

\begin{lstlisting}[style=shellcode, caption={Outputs surfaced by the root module after apply}]
$ terraform output
lb_public_ip        = "104.211.49.157"
resource_group_name = "lab8-rg"
vm_names            = [
  "lab8-vm-0",
  "lab8-vm-1",
]
\end{lstlisting}

A small loop hitting the Load Balancer six times demonstrates that
successive requests rotate between both backends, confirming the
backend pool is healthy and the 5-tuple hash distributes flows as
expected.

\begin{lstlisting}[style=shellcode, caption={Round-robin verification against the LB Public IP}]
$ for ($i=0; $i -lt 6; $i++) { curl http://104.211.49.157 }
<h1>Hello from lab8-vm-0</h1>
<h1>Hello from lab8-vm-1</h1>
<h1>Hello from lab8-vm-0</h1>
<h1>Hello from lab8-vm-1</h1>
<h1>Hello from lab8-vm-0</h1>
<h1>Hello from lab8-vm-1</h1>
\end{lstlisting}

% ---------------------------------------------------------------------------
\section{Security Posture}
% ---------------------------------------------------------------------------

Even at lab scale, the team applied a baseline that mirrors what would
be required in production:

\begin{itemize}
  \item \textbf{NSG with explicit allow-list.} Only 80/TCP from the
        public Internet (required for the LB) and 22/TCP from the
        operator's \texttt{/32} are permitted; everything else is
        denied by default.
  \item \textbf{SSH key-only authentication.} Password authentication
        is disabled by the \texttt{azurerm} provider when
        \texttt{admin\_ssh\_key} is supplied, structurally preventing
        password spray attacks.
  \item \textbf{Tagged ephemeral resources.} Every resource carries
        \texttt{owner}, \texttt{course}, \texttt{env} and
        \texttt{expires} tags, making cost ownership and cleanup
        auditable.
  \item \textbf{No long-lived secret in CI/CD.} The OIDC federation
        means there is no Service Principal client secret to rotate or
        leak.
  \item \textbf{NIC-level NSG attachment.} Avoids subnet-level
        inheritance pitfalls when adding new tiers.
\end{itemize}

\begin{warningbox}
Exposing 22/TCP to the Internet, even restricted to a single \texttt{/32},
is still an attack surface. For production, the lab leaves
\texttt{subnet-mgmt} pre-provisioned so \textbf{Azure Bastion} can be
introduced later, removing public SSH entirely.
\end{warningbox}

% ---------------------------------------------------------------------------
\section{Cost Analysis}
% ---------------------------------------------------------------------------

Approximate East US list prices, validated against the Azure pricing
calculator on the deployment date:

\begin{center}
\small
\begin{tabular}{@{}lrrr@{}}
\toprule
\textbf{Resource} & \textbf{Hourly} & \textbf{Daily} & \textbf{Monthly} \\
\midrule
2 $\times$ VM \texttt{Standard\_B1s} & \$0.0104 & \$0.50 & \$15.20 \\
Standard Load Balancer            & \$0.0250 & \$0.60 & \$18.30 \\
Standard Public IP (static)       & \$0.0050 & \$0.12 & \$3.70 \\
Storage Account (state, LRS)      & \$0.0002 & \$0.005 & \$0.15 \\
\midrule
\textbf{Total}                    & \textbf{\$0.041} & \textbf{\$1.23} & \textbf{\$37.35} \\
\bottomrule
\end{tabular}
\end{center}

An eight-hour development session therefore costs roughly USD \$0.33,
which fits comfortably inside the USD \$100 Azure for Students credit
and allows dozens of full lab cycles before the credit is exhausted.

% ---------------------------------------------------------------------------
\section{Optional Challenges Implemented}
% ---------------------------------------------------------------------------

The lab guide lists five optional challenges (\emph{Retos}) and asks
students to pick two or more. The team implemented two of them as
\textbf{opt-in modules behind feature flags}, so the cheap base lab
remains untouched and the additional resources only spin up when the
team explicitly enables them. The full enable/disable procedure is
documented in \texttt{docs/CHALLENGES.md}.

\subsection{Challenge \#1 --- Azure Bastion (no public SSH)}

The first challenge eliminates the public 22/TCP attack surface on the
workload VMs. When \texttt{enable\_bastion = true}, the \texttt{vnet}
module provisions an additional subnet named exactly
\texttt{AzureBastionSubnet} (\texttt{10.10.3.0/26}, both name and
minimum size mandated by Azure), and a new \texttt{bastion} module
deploys a Standard SKU Public IP plus the
\texttt{azurerm\_bastion\_host} itself. Operators reach the VMs by
clicking \emph{Connect $\rightarrow$ Bastion} in the Azure portal, which
brokers SSH/RDP over TLS without ever exposing port 22 to the Internet.

\begin{lstlisting}[style=hclcode, caption={Bastion module wired through a feature flag (\texttt{infra/main.tf})}]
module "bastion" {
  count               = var.enable_bastion ? 1 : 0
  source              = "../modules/bastion"
  resource_group_name = azurerm_resource_group.rg.name
  location            = azurerm_resource_group.rg.location
  prefix              = var.prefix
  subnet_id           = module.vnet.subnet_bastion_id
  sku                 = var.bastion_sku
  tags                = var.tags
}
\end{lstlisting}

\begin{warningbox}
Azure Bastion is the most expensive single resource in this lab
(approximately USD \$0.19/hour for the Basic SKU). The team kept the
flag disabled by default so a forgotten \texttt{terraform apply} cannot
silently drain the student credit; production deployments would also
remove the \texttt{Allow-SSH-From-Operator} NSG rule entirely once
Bastion is in place.
\end{warningbox}

\subsection{Challenge \#2 --- Cost Management budget with email alerts}

The second challenge installs a hard \emph{guardrail against runaway
spend}: a Cost Management budget scoped to the lab Resource Group with
three notification thresholds (50\,\%, 90\,\% and 100\,\% of the
configured monthly cap), all evaluated on \emph{Actual} cost rather
than \emph{Forecast} to avoid noisy notifications during the first few
days of the month. Budgets themselves are free; the only added cost is
the email pipeline (also free).

\begin{lstlisting}[style=hclcode, caption={Budget module: three notifications, gated by feature flag}]
module "budget" {
  count             = var.enable_budget ? 1 : 0
  source            = "../modules/budget"
  prefix            = var.prefix
  resource_group_id = azurerm_resource_group.rg.id
  amount_usd        = var.budget_amount_usd
  contact_emails    = var.budget_contact_emails
  start_date        = var.budget_start_date
}

# inside modules/budget/main.tf:
notification {
  enabled        = true
  threshold      = 50
  operator       = "GreaterThan"
  threshold_type = "Actual"
  contact_emails = var.contact_emails
}
\end{lstlisting}

\begin{resultbox}[title={Why these two challenges?}]
\begin{itemize}
  \item \textbf{Bastion} attacks the largest residual security risk in
        the base architecture (the public SSH surface) without
        introducing a new technology stack to the team.
  \item \textbf{Budget alert} adds operational maturity at zero
        marginal cost --- the perfect ergonomics win for a lab whose
        primary financial constraint is a fixed student credit.
  \item Both are wired with \texttt{count = var.enable\_X ? 1 : 0},
        which keeps the \texttt{terraform plan} for the base lab
        byte-identical to the pre-challenge version when the flags
        are off.
\end{itemize}
\end{resultbox}

% ---------------------------------------------------------------------------
\section{Video Evidence}
% ---------------------------------------------------------------------------

A full demo recording captures the end-to-end execution: backend
bootstrap, \texttt{terraform init} against the remote backend,
\texttt{terraform apply}, \texttt{curl} loop verifying round-robin and
\texttt{terraform destroy} for the cleanup. The video is hosted on
Google Drive and linked from the repository \texttt{README}:

\begin{resultbox}[title={Demo recording}]
\textbf{URL:} \url{https://drive.google.com/file/d/1CAY1gmPZbe96VobD-xGRg0q5HYMfCjOu/view}
\medskip

\noindent The recording showcases:
\begin{itemize}
  \item Storage Account / container provisioning for remote state.
  \item Terraform initialising the AzureRM backend and downloading
        provider \texttt{azurerm} v4.72.0.
  \item Successful \texttt{plan} and \texttt{apply} producing the
        Resource Group, VNet, NICs, VMs, LB and NSG.
  \item Live \texttt{curl} loop alternating between
        \texttt{lab8-vm-0} and \texttt{lab8-vm-1}.
  \item Clean teardown via \texttt{terraform destroy} confirming a
        zero-leftover state.
\end{itemize}
\end{resultbox}

% ---------------------------------------------------------------------------
\section{Conclusions}
% ---------------------------------------------------------------------------

\begin{resultbox}[title={Key Findings}]
\begin{enumerate}
  \item \textbf{IaC dramatically reduces operational risk.} Replacing
        portal click-ops with declarative Terraform turned an
        error-prone, multi-step procedure into a single
        \texttt{apply}, fully reproducible and reviewable as code.
        Destroying the deployment is just as deterministic, which is
        critical for cost control on student credits.
  \item \textbf{Module decomposition pays off immediately.} Splitting
        the workload into \texttt{vnet}, \texttt{compute} and
        \texttt{lb} kept each diff small and surfaced the precise
        seam at which a future Application Gateway upgrade would land.
  \item \textbf{OIDC is the right default for cloud CI/CD.} Removing
        the long-lived Service Principal secret eliminated an entire
        class of credential-management problems. Federated credentials
        scoped per branch and per environment make the trust boundary
        easy to reason about.
  \item \textbf{Standard L4 LB is the right tool for this workload.}
        Application Gateway would have multiplied cost and operational
        surface for routing features we do not consume. The clean
        module boundary keeps the upgrade path open.
\end{enumerate}
\end{resultbox}

\subsection{Lessons Learned}

\begin{itemize}
  \item Pinning the AzureRM provider via \texttt{.terraform.lock.hcl}
        is non-negotiable: a minor provider bump silently changed
        defaults during early experimentation and caused unexpected
        plan diffs.
  \item NIC-level NSG attachment is more verbose than subnet-level,
        but completely avoids the NSG inheritance debugging that
        wasted time during earlier exercises.
  \item \texttt{terraform fmt -recursive} should be a pre-commit hook
        from day one; the CI \texttt{fmt -check} block in the Actions
        workflow catches drift but the developer feedback loop is
        much faster locally.
  \item Validating variables with \texttt{validation} blocks turns the
        \texttt{plan} step into a spec test: bad input fails before a
        single Azure API call is issued, which is exactly the
        feedback loop we want.
\end{itemize}

% ---------------------------------------------------------------------------
\section{References}
% ---------------------------------------------------------------------------

\begin{enumerate}
  \item HashiCorp. \emph{Terraform Language Documentation}.
        \url{https://developer.hashicorp.com/terraform/language}

  \item Microsoft Azure. \emph{Azure Load Balancer overview}.
        \url{https://learn.microsoft.com/en-us/azure/load-balancer/load-balancer-overview}

  \item Microsoft Azure. \emph{Use Azure Storage as a Terraform backend}.
        \url{https://learn.microsoft.com/en-us/azure/developer/terraform/store-state-in-azure-storage}

  \item Microsoft Azure. \emph{Configure a federated identity credential
        on an app for GitHub Actions}.
        \url{https://learn.microsoft.com/en-us/azure/developer/github/connect-from-azure-openid-connect}

  \item HashiCorp. \emph{AzureRM Provider Documentation (v4)}.
        \url{https://registry.terraform.io/providers/hashicorp/azurerm/latest/docs}

  \item Canonical. \emph{cloud-init Documentation}.
        \url{https://cloudinit.readthedocs.io/en/latest/}

  \item GitHub Docs. \emph{About security hardening with OpenID Connect}.
        \url{https://docs.github.com/en/actions/deployment/security-hardening-your-deployments/about-security-hardening-with-openid-connect}
\end{enumerate}

\end{document}
